\chapter{错误处理与传播}
\label{ch:16}

缺少一个值和执行失败是不同状态。Option 表达有无，Result 表达成功或错误；接口错误让调用边界统一返回类型，同时保留具体错误供恢复分支识别。

\section{Option 与 Result 的分支}

默认构造的 Option 从其默认变体和值开始，不应当作显式缺省；需要缺省时写 Option<T>::None。读取可选值前匹配 Value／None，读取结果前匹配 Ok／Err。unwrap 适用于已经建立前提的位置，不应代替外部输入校验。下面把缺省处理与逐层返回组合。

\par
\medskip
\begin{lstlisting}[style=chali]
Option<int64> findFirst(List<int64> xs, int64 target)
{
    for (int64 x in xs)
    {
        if (x == target)
        {
            return Value(x);
        }
    }
    return None;
}

void main()
{
    List<int64> xs;
    xs.push(3);
    xs.push(7);
    println("{}", findFirst(xs, 7).unwrap());
    println("{}", findFirst(xs, 9).unwrapOr(-1));
    match (findFirst(xs, 3))
    {
        Value(v) => { println("found = {}", v); }
        None => { println("no"); }
    }
}
\end{lstlisting}
\par
\medskip

\par
\medskip
\begin{lstlisting}[style=out]
7
-1
found = 3
\end{lstlisting}
\par
\medskip

\section{枚举载荷中的接口匹配}

对于 Result<String, Error>，普通成功载荷直接写 Ok(text)。错误载荷需要区分实现时，可以写 Err(IoError e)，单载荷变体也支持 Err::IoError(e)。在可唯一解析、不与变体或其他候选冲突时，还可以省略外层写 IoError(e)。剩余错误用 Err(\_) 或 Err(\_ other) 接收。

这些省略语法仅用于单载荷变体；多个载荷必须保留外层变体和括号内类型标注。不能凭一个类型名猜测它对应多个载荷中的哪一项。

\par
\medskip
\begin{lstlisting}[style=chali]
import std::io;
Result<String, Error> readDemo()
{
    return Result<String, Error>::Err(std::io::IoError::ReadFailed("demo"));
}
void main()
{
    let result = readDemo();
    match (result)
    {
        Ok(text) => { println("{}", text); }
        Err::std::io::IoError(e) => { println("qualified: {}", e.toString()); }
        Err(_) => { println("other error"); }
    }
    match (result)
    {
        Ok(_) => { println("ok"); }
        std::io::IoError(e) => { println("short: {}", e.toString()); }
        Err(_) => { println("other error"); }
    }
}
\end{lstlisting}
\par
\medskip

\par
\medskip
\begin{lstlisting}[style=out]
qualified: 读取失败：demo
short: 读取失败：demo
\end{lstlisting}
\par
\medskip

两个载荷保留外层 Failed(code, Type value)。当 Ok 和 Err 都可能带同一个接口实现时，显式变体区分标签，避免省略后出现歧义。

\par
\medskip
\begin{lstlisting}[style=chali]
import std::io;
enum Response
{
    Pending,
    Failed(int64 code, Error error)
};
void main()
{
    let response = Response::Failed(503, std::io::IoError::Closed);
    match (response)
    {
        Response::Pending => { println("pending"); }
        Response::Failed(code, std::io::IoError e) => { println("{}: {}", code, e.toString()); }
        Response::Failed(code, _ other) => { println("{}: {}", code, other.toString()); }
    }
    Result<Error, Error> result = Result<Error, Error>::Ok(std::io::IoError::Closed);
    match (result)
    {
        Ok(std::io::IoError e) => { println("ok: {}", e.toString()); }
        Ok(_) => { println("other success"); }
        Err(std::io::IoError e) => { println("err: {}", e.toString()); }
        Err(_) => { println("other failure"); }
    }
}
\end{lstlisting}
\par
\medskip

\par
\medskip
\begin{lstlisting}[style=out]
503: 文件已关闭
ok: 文件已关闭
\end{lstlisting}
\par
\medskip

\section{传播时的接口提升}

? 在成功路径取得内部值，在失败或缺省路径按传播协议提前返回。若内部错误 IoError 实现了返回位置所需的 Error，Result<..., IoError> 可以向 Result<..., Error> 传播。Option 的接口载荷和自定义传播协议同样需要满足目标类型及协议约束，不能借 ? 跳过类型检查。

leaf 的具体 IoError 通过 ? 提升为外层 Error。Option 的成功载荷也能转为接口值；缺省路径仍由 ? 的 Option 协议处理。

\par
\medskip
\begin{lstlisting}[style=chali]
import std::io;
Result<String, std::io::IoError> leaf()
{
    return Result<String, std::io::IoError>::Err(std::io::IoError::Closed);
}
Result<String, Error> readName()
{
    String name = leaf()?;
    return Result<String, Error>::Ok(name);
}
Option<std::io::IoError> optionalError()
{
    return Option<std::io::IoError>::Value(std::io::IoError::Closed);
}
Option<Error> widenedError()
{
    let error = optionalError()?;
    return Option<Error>::Value(error);
}
void main()
{
    match (readName())
    {
        Ok(name) => { println("{}", name); }
        Err(std::io::IoError e) => { println("result: {}", e.toString()); }
        Err(_) => { println("other error"); }
    }
    match (widenedError())
    {
        Value(std::io::IoError e) => { println("option: {}", e.toString()); }
        Value(_) => { println("other value"); }
        None => { println("absent"); }
    }
}
\end{lstlisting}
\par
\medskip

\par
\medskip
\begin{lstlisting}[style=out]
result: 文件已关闭
option: 文件已关闭
\end{lstlisting}
\par
\medskip

\section{自定义传播与恢复边界}

自定义 Try 协议需要同时定义成功结果与残余状态如何转换。把恢复放在能做出决定的位置：文件错误可以重试或显示路径，业务错误可以返回字段提示，未知实现保留统一错误显示。清理规则继续适用于提前返回。

\par
\medskip
\begin{lstlisting}[style=chali]
enum Step: Try<int64, NoneResidual>
{
    Value(int64 value), Empty;
    TryBranch<int64, NoneResidual> branch()
    {
        match (self)
        {
            Step::Value(value) =>
            {
                return TryBranch<int64, NoneResidual>::Continue(value);
            }
            Step::Empty =>
            {
                return TryBranch<int64, NoneResidual>::Break(NoneResidual::None);
            }
        }
        panic("unreachable");
    }
};
Option<int64> addOne(Step step)
{
    int64 value = step?;
    return Option<int64>::Value(value + 1);
}
void main()
{
    println("{} {}", addOne(Step::Value(4)).unwrap(), addOne(Step::Empty).isNone());
}
\end{lstlisting}
\par
\medskip

\par
\medskip
\begin{lstlisting}[style=out]
5 1
\end{lstlisting}
\par
\medskip

\section{习题}

\begin{enumerate}

\item 把 Err(\_) 改成带名字的回退。

\item 何时能省略外层 Err？

\item 两个载荷的错误分支应如何标注接口类型？

\item ? 是否能把无关错误类型自动当成 Error？

\end{enumerate}

\subsection*{参考解答与验收要点}

\begin{enumerate}

\item 使用 Err(\_ other)，再调用 other.toString()。

\item 单载荷且解析唯一，没有变体或类型候选冲突。

\item 保留变体，例如 Failure(code, IoError error)。

\item 不能，必须实现对应接口并满足传播协议。

\end{enumerate}

下一章：\hyperref[ch:17]{第 17 章}。

\section{番外：错误又来了}

改写自伊索寓言《狼来了》。

一个放羊的孩子在山坡上看羊，觉得空闲无聊，便向村里大喊：“成功了！成功了！”村人赶来，才发现网络已经断了，返回值只是孩子随手填的零。

第二天，孩子又喊成功。村人问查询结果，孩子答：“零只羊。”问是不是空集合，答：“也许。”问是不是没有对象，答：“也许。”问是不是操作失败，孩子拍着终端道：“反正没有错误，错误都被我改成零了。”

村人渐渐不再看他的报信。到了狼真来时，警报查询恰好失败。孩子照例送回零，仓库照例撤销警报。羊群照例遭殃，只有监控图表一片安宁。

老村长气得胡子直抖，给孩子订了规矩：成功就带数，未找到就说未找到，失败就带具体错误。别让 \texttt{Option}、\texttt{Result} 和零值三家挤在一个门牌下。

孩子问：“那么传给上一层之后，狼是不是就走了？”村长说：“\texttt{?} 是传话，不是驱狼。上层还得决定重试、求援还是撤离。”

后来村里每个错误都实现了同一个 \texttt{Error} 接口，却仍保留各自的细节。狼又来时，孩子终于能说清是哪一路失败。村长颇欣慰，只叮嘱一句：“别为了报表好看，又给我统一成‘未知异常’。”

\subsection*{技术对照}

Option 表达存在与否，Result 表达成功或失败。错误传播不等于已经恢复；unwrap、assert 和 panic 适合的前提需要由调用位置明确保证。
